\documentclass[11pt]{article}
\usepackage{amsmath,amssymb,amsthm}
\usepackage{algorithm}
\usepackage[noend]{algpseudocode}
\usepackage{enumitem}
\usepackage{hyperref}
\usepackage{geometry}
\geometry{margin=1in}
\newtheorem{theorem}{Theorem}
\newtheorem{lemma}{Lemma}
\newtheorem{remark}{Remark}
\newcommand{\R}{\mathbb{R}}
\newcommand{\Cset}{\mathcal{C}}
\newcommand{\gap}{\mathcal{G}}
\begin{document}

\section*{Frank–Wolfe (Conditional Gradient) Method}

\section*{Mathematical Formulation}
Let $f:\R^{d}\rightarrow\R$ be a convex and $L$–smooth function and let 
\[
\Cset\subset\R^{d}
\]
be a non‑empty compact convex set (the feasible domain).  
At iteration $k\ge 0$ we have the current point $x_{k}\in\Cset$.  
The algorithm performs the following steps:

\begin{enumerate}[label=\arabic*)]
\item \textbf{Linear minimization oracle (LMO)}  
\[
s_{k}\;:=\;\arg\min_{s\in\Cset}\;\langle \nabla f(x_{k}),\,s\rangle .
\tag{1}
\]
\item \textbf{Step–size selection} $\gamma_{k}\in[0,1]$ (e.g. $\gamma_{k}=2/(k+2)$ or exact line search).  
\item \textbf{Update}  
\[
x_{k+1}\;:=\;(1-\gamma_{k})\,x_{k}+\gamma_{k}\,s_{k}.
\tag{2}
\]
\end{enumerate}

The algorithm never leaves $\Cset$ because $x_{k+1}$ is a convex combination of two points in $\Cset$.

\section*{Pseudocode}
\begin{algorithm}[H]
\caption{Frank–Wolfe (Conditional Gradient)}\label{alg:FW}
\begin{algorithmic}[1]
\Require Convex $L$‑smooth $f$, compact convex $\Cset$, initial $x_{0}\in\Cset$, 
step–size rule $\{\gamma_{k}\}_{k\ge0}$.
\For{$k=0,1,2,\dots$}
    \State Compute gradient $g_{k}\gets\nabla f(x_{k})$.
    \State \textbf{LMO:} $s_{k}\gets\arg\min_{s\in\Cset}\langle g_{k},s\rangle$.
    \State Choose step size $\gamma_{k}\in[0,1]$.
    \State Update $x_{k+1}\gets (1-\gamma_{k})x_{k}+\gamma_{k}s_{k}$.
    \If{stopping criterion satisfied (e.g. duality gap $<\varepsilon$)}
        \State \textbf{break}
    \EndIf
\EndFor
\State \Return $x_{k}$.
\end{algorithmic}
\end{algorithm}

\section*{Dry Run Example}
Consider the one–dimensional problem  

\[
f(w)=(w-3)^{2},\qquad \Cset=[0,5],\qquad w_{0}=0,\qquad\gamma_{k}=0.1\;\;(\forall k).
\]

\begin{center}
\begin{tabular}{c|c|c|c|c}
$k$ & $w_{k}$ & $\nabla f(w_{k})=2(w_{k}-3)$ & $s_{k}=\arg\min_{s\in[0,5]}\langle\nabla f(w_{k}),s\rangle$ & $w_{k+1}=(1-\gamma)w_{k}+\gamma s_{k}$ \\ \hline
0 & $0$ & $-6$ & $5$ (since $-6s$ is minimal at $s=5$) & $0.5$ \\
1 & $0.5$ & $-5$ & $5$ & $0.95$ \\
2 & $0.95$ & $-4.1$ & $5$ & $1.355$ \\
\end{tabular}
\end{center}

Corresponding objective values:

\[
f(0)=9,\quad f(0.5)=6.25,\quad f(0.95)=4.2025,\quad f(1.355)=2.704\;.
\]

The iterates move towards the minimizer $w^{*}=3$ while remaining in the feasible interval.

\newpage
\section*{Assumptions}
\begin{enumerate}[label=\textbf{A\arabic*:}, leftmargin=*]
\item \label{A1} \textbf{Convexity.} $f$ is convex on $\Cset$:
\[
f(y)\ge f(x)+\langle\nabla f(x),y-x\rangle,\qquad\forall x,y\in\Cset.
\]
\item \label{A2} \textbf{Smoothness.} $f$ has $L$–Lipschitz gradient on $\Cset$:
\[
\|\nabla f(y)-\nabla f(x)\| \le L\|y-x\|,\qquad\forall x,y\in\Cset.
\]
Equivalently (Descent Lemma):
\[
f(y)\le f(x)+\langle\nabla f(x),y-x\rangle+\frac{L}{2}\|y-x\|^{2}.
\tag{3}
\]
\item \label{A3} \textbf{Compactness.} $\Cset$ is compact; denote its diameter
\[
D:=\max_{x,y\in\Cset}\|x-y\|<\infty .
\]
\item \label{A4} \textbf{Existence of optimal solution.} There exists $x^{*}\in\Cset$ such that
\[
f(x^{*})=\min_{x\in\Cset}f(x).
\]
\end{enumerate}

\section*{Main Convergence Theorem}
\begin{theorem}[Standard Frank–Wolfe rate]\label{thm:FW}
Under Assumptions \ref{A1}–\ref{A4} define the curvature constant
\[
C_{f}:=\sup_{\substack{x,s\in\Cset\\ \gamma\in[0,1]}}
\frac{2}{\gamma^{2}}\Bigl( f\bigl(x+\gamma(s-x)\bigr)-f(x)-\gamma\langle\nabla f(x),s-x\rangle\Bigr).
\tag{4}
\]
If the step size is chosen as $\displaystyle\gamma_{k}= \frac{2}{k+2}$, then for all $k\ge0$
\[
f(x_{k})-f(x^{*})\le \frac{2C_{f}}{k+2}.
\tag{5}
\]
Consequently $f(x_{k})\to f(x^{*})$ at a rate $\mathcal{O}(1/k)$.
\end{theorem}

\section*{Theoretical Properties}
\begin{enumerate}[label=\textbf{P\arabic*:}, leftmargin=*]
\item \textbf{Primal–dual gap.} The Frank–Wolfe (duality) gap
\[
\gap(x_{k})\;:=\;\langle \nabla f(x_{k}),x_{k}-s_{k}\rangle
\tag{6}
\]
satisfies $0\le\gap(x_{k})\le f(x_{k})-f(x^{*})$, thus can be used as a stopping
criterion.
\item \textbf{Stability w.r.t. step size.} Any step size sequence
$\gamma_{k}\in[0,1]$ yields a non‑increasing objective; the specific choice
$\gamma_{k}=2/(k+2)$ guarantees the optimal $\mathcal{O}(1/k)$ bound.
\item \textbf{Comparison to projected gradient.} Frank–Wolfe avoids
projections onto $\Cset$, requiring only a linear minimization oracle. The
price is the slower $\mathcal{O}(1/k)$ rate compared with $\mathcal{O}(1/k)$ for
projected gradient under the same smoothness, but with cheaper per‑iteration
cost when the LMO is simple.
\item \textbf{Extension.} With strong convexity of $f$ and a polyhedral
$\Cset$, the rate can be accelerated to $\mathcal{O}(1/k^{2})$ using
away‑steps or pairwise variants.
\end{enumerate}

\section*{Discussion}
The theorem hinges on the curvature constant $C_{f}$, which is bounded by
$L D^{2}$ (by Lemma~\ref{lem:curv-bound}). Hence (5) can be rewritten as
$f(x_{k})-f(x^{*})\le \frac{2LD^{2}}{k+2}$. The bound is independent of the
dimension and holds for any compact convex feasible set as long as the LMO can
be solved exactly.

\newpage
\section*{Preliminaries (Definitions and Lemmas)}
\begin{lemma}[Descent Lemma]\label{lem:descent}
Under Assumption \ref{A2}, for any $x,y\in\Cset$,
\[
f(y)\le f(x)+\langle\nabla f(x),y-x\rangle+\frac{L}{2}\|y-x\|^{2}.
\]
\end{lemma}
\begin{lemma}[Curvature bound]\label{lem:curv-bound}
For $L$‑smooth $f$ on a set of diameter $D$, the curvature constant satisfies
\[
C_{f}\le L D^{2}.
\]
\end{lemma}
\begin{proof}
Let $x,s\in\Cset$ and $\gamma\in[0,1]$. Using Lemma~\ref{lem:descent} with
$y=x+\gamma(s-x)$ gives
\[
f(x+\gamma(s-x))\le f(x)+\gamma\langle\nabla f(x),s-x\rangle+\frac{L\gamma^{2}}{2}\|s-x\|^{2}.
\]
Rearranging yields
\[
\frac{2}{\gamma^{2}}\bigl(f(x+\gamma(s-x))-f(x)-\gamma\langle\nabla f(x),s-x\rangle\bigr)
\le L\|s-x\|^{2}\le LD^{2}.
\]
Taking the supremum over $x,s,\gamma$ proves the claim. \hfill\qedsymbol
\end{proof}

\section*{Main Proof (Step‑by‑Step Derivation)}
We prove Theorem~\ref{thm:FW}. Throughout the proof we fix $k\ge0$.

\begin{enumerate}[label=[Step \arabic*], leftmargin=*]
\item \textbf{Apply smoothness to the update point.}  
Using Lemma~\ref{lem:descent} with $x=x_{k}$ and $y=x_{k+1}=x_{k}+\gamma_{k}(s_{k}-x_{k})$,
\[
f(x_{k+1})\le f(x_{k})+\gamma_{k}\langle\nabla f(x_{k}),s_{k}-x_{k}\rangle
+\frac{L\gamma_{k}^{2}}{2}\|s_{k}-x_{k}\|^{2}.
\tag{7}
\]

\item \textbf{Replace the quadratic term by curvature.}  
By definition of $C_{f}$ (4),
\[
\frac{L}{2}\|s_{k}-x_{k}\|^{2}\le \frac{C_{f}}{2},
\]
hence
\[
\frac{L\gamma_{k}^{2}}{2}\|s_{k}-x_{k}\|^{2}\le \frac{C_{f}\gamma_{k}^{2}}{2}.
\tag{8}
\]

\item \textbf{Use optimality of the LMO.}  
From (1) we have
\[
\langle\nabla f(x_{k}),s_{k}\rangle
= \min_{s\in\Cset}\langle\nabla f(x_{k}),s\rangle
\le \langle\nabla f(x_{k}),x^{*}\rangle .
\tag{9}
\]
Rearranging,
\[
\langle\nabla f(x_{k}),s_{k}-x_{k}\rangle
\le \langle\nabla f(x_{k}),x^{*}-x_{k}\rangle .
\tag{10}
\]

\item \textbf{Convexity inequality.}  
By convexity (Assumption \ref{A1}),
\[
f(x^{*})\ge f(x_{k})+\langle\nabla f(x_{k}),x^{*}-x_{k}\rangle .
\tag{11}
\]
Thus
\[
\langle\nabla f(x_{k}),x^{*}-x_{k}\rangle\le f(x^{*})-f(x_{k}).
\tag{12}
\]

\item \textbf{Combine (7)–(12).}  
Insert (10) and (12) into (7) and then (8):
\[
\begin{aligned}
f(x_{k+1})
&\le f(x_{k})+\gamma_{k}\bigl(f(x^{*})-f(x_{k})\bigr)
+\frac{C_{f}\gamma_{k}^{2}}{2}\\[4pt]
&= (1-\gamma_{k})f(x_{k})+\gamma_{k}f(x^{*})+\frac{C_{f}\gamma_{k}^{2}}{2}.
\end{aligned}
\tag{13}
\]

\item \textbf{Subtract $f(x^{*})$ and define the error.}  
Let $h_{k}:=f(x_{k})-f(x^{*})\ge0$.  From (13),
\[
h_{k+1}\le (1-\gamma_{k})h_{k}+ \frac{C_{f}\gamma_{k}^{2}}{2}.
\tag{14}
\]

\item \textbf{Plug the specific step size $\gamma_{k}=2/(k+2)$.}  
Compute
\[
1-\gamma_{k}=1-\frac{2}{k+2}= \frac{k}{k+2},\qquad
\gamma_{k}^{2}= \frac{4}{(k+2)^{2}}.
\tag{15}
\]

\item \textbf{Recurrence simplification.}  
Insert (15) into (14):
\[
h_{k+1}\le \frac{k}{k+2}\,h_{k}+ \frac{C_{f}}{2}\cdot\frac{4}{(k+2)^{2}}
= \frac{k}{k+2}\,h_{k}+ \frac{2C_{f}}{(k+2)^{2}}.
\tag{16}
\]

\item \textbf{Induction to obtain closed form.}  
We prove by induction that
\[
h_{k}\le \frac{2C_{f}}{k+2},\qquad\forall k\ge0.
\tag{17}
\]

\begin{itemize}
\item \textbf{Base case $k=0$.}  
From (14) with $\gamma_{0}=1$ (since $2/(0+2)=1$) we have
\[
h_{1}\le (1-1)h_{0}+ \frac{C_{f}}{2}\cdot1 = \frac{C_{f}}{2}
\le \frac{2C_{f}}{2}= \frac{2C_{f}}{1+2},
\]
so (17) holds for $k=1$. Because $h_{0}\le C_{f}$ (by definition of curvature) the bound also holds for $k=0$.

\item \textbf{Inductive step.}  
Assume (17) holds for some $k\ge0$. Using (16),
\[
\begin{aligned}
h_{k+1}
&\le \frac{k}{k+2}\cdot\frac{2C_{f}}{k+2}+ \frac{2C_{f}}{(k+2)^{2}}\\[4pt]
&= \frac{2C_{f}k}{(k+2)^{2}}+\frac{2C_{f}}{(k+2)^{2}}
= \frac{2C_{f}(k+1)}{(k+2)^{2}}\\[4pt]
&= \frac{2C_{f}}{(k+1)+2},
\end{aligned}
\]
which is exactly the statement (17) for index $k+1$. Hence the induction
closes.
\end{itemize}
Thus (17) is proved, which is precisely inequality (5). \hfill\qedsymbol
\end{enumerate}

\section*{Rate Derivation (With Constants)}
From Lemma~\ref{lem:curv-bound} we have $C_{f}\le LD^{2}$. Substituting into (5) yields
\[
f(x_{k})-f(x^{*})\le \frac{2LD^{2}}{k+2},
\qquad k\ge0.
\tag{18}
\]
Therefore the Frank–Wolfe method achieves a sublinear convergence rate
$\mathcal{O}\!\bigl(\tfrac{1}{k}\bigr)$ with an explicit constant $2LD^{2}$.

\section*{Special Cases}
\begin{enumerate}[label=\textbf{S\arabic*:}, leftmargin=*]
\item \textbf{Exact line search.}  
If $\gamma_{k}$ is chosen as the minimizer of $f\bigl((1-\gamma)x_{k}+\gamma s_{k}\bigr)$ over $\gamma\in[0,1]$, the same recurrence (14) holds with a possibly smaller $\gamma_{k}$, thus the bound (5) still applies.
\item \textbf{Polyhedral $\Cset$ and strong convexity.}  
When $f$ is $\mu$‑strongly convex and $\Cset$ is a polytope, the away‑step Frank–Wolfe algorithm enjoys a linear rate
\[
f(x_{k})-f(x^{*})\le \bigl(1-\rho\bigr)^{k}\bigl(f(x_{0})-f(x^{*})\bigr),
\]
with $\rho>0$ depending on $\mu$, $L$, and the pyramidal width of $\Cset$.
\item \textbf{Stochastic Frank–Wolfe.}  
If the gradient is replaced by an unbiased stochastic estimator with bounded variance, the same analysis yields an expected rate
\[
\mathbb{E}\bigl[f(\bar{x}_{T})-f(x^{*})\bigr]=\mathcal{O}\!\bigl(\tfrac{1}{\sqrt{T}}\bigr),
\]
mirroring the classic SGD result.
\end{enumerate}

\end{document}